\section*{Hoofdstuk \ref{ch:b2:coordination-complexes} --- Overgangsmetalen en coördinatieverbindingen}

\begin{solution}{exo:b2:coordination-complexes:1}
\ce{Ti^3+} $d^1$, \ce{V^3+} $d^2$, \ce{Cr^3+} $d^3$, \ce{Mn^2+} $d^5$, \ce{Fe^3+} $d^5$, \ce{Co^2+}
$d^7$, \ce{Ni^2+} $d^8$, \ce{Zn^2+} $d^{10}$.
\end{solution}

\begin{solution}{exo:b2:coordination-complexes:2}
Water 1, en 2, ethaandioaat 2, edta 6, thiocyanaat 1; thiocyanaat is ambidentaat
(via S of N).
\end{solution}

\begin{solution}{exo:b2:coordination-complexes:3}
Tetraamminekoper(II)sulfaat; kalium\-hexa\-cyanido\-ferraat(III);
tetraammine-dichlorido\-kobalt(III)\-chloride; tetracarbonylnikkel(0).
\end{solution}

\begin{solution}{exo:b2:coordination-complexes:4}
Lineair; tetraëdrisch; vlak vierkant ($d^8$, derde reeks); octaëdrisch.
\end{solution}

\begin{solution}{exo:b2:coordination-complexes:5}
Twee: cis (de twee chloriden naast elkaar) en trans (tegenover elkaar). In een tetraëder zou
er maar één zijn.
\end{solution}

\begin{solution}{exo:b2:coordination-complexes:6}
\ce{Cr(CO)6} 18; \ce{Fe(CO)5} 18; \ce{Ni(CO)4} 18; \ce{V(CO)6}: $5 + 12 = 17$;
\ce{Mn2(CO)10} 18 per mangaanatoom. \ce{V(CO)6} overtreedt de regel: het is \omterm{def:b2:diatomic-mos:magnetism}{paramagnetisch} en
wordt gemakkelijk gereduceerd tot \ce{[V(CO)6]-}, dat er 18 heeft.
\end{solution}

\begin{solution}{exo:b2:coordination-complexes:7}
Ferroceen: Fe(II) $d^6$ + 12 = 18. Kobaltoceen: Co(II) $d^7$ + 12 = 19. Kobaltoceen verliest
zijn extra elektron gemakkelijk en geeft het kobaltoceniumion \ce{[Co(C5H5)2]+}, met 18.
\end{solution}

\begin{solution}{exo:b2:coordination-complexes:8}
Eén edta vervangt zes watermoleculen: \ce{[Ni(H2O)6]^2+ + edta^4- -> [Ni(edta)]^2- +
6H2O}, zeven deeltjes uit twee. De \omterm{def:b2:reaction-free-energy:reaction-entropy}{reactie-entropie} is sterk positief, en
$K$ zeer groot. Zes afzonderlijke liganden vervangen zes watermoleculen zonder winst in het
aantal deeltjes.
\end{solution}

\begin{solution}{exo:b2:coordination-complexes:9}
Nitrito-$\kappa N$ (\ce{Co-NO2}, het nitro-isomeer, geel) en nitrito-$\kappa O$
(\ce{Co-O-N=O}, het nitrito-isomeer, rood).
\end{solution}

\begin{solution}{exo:b2:coordination-complexes:10}
Drie: het trans-isomeer (achiraal: spiegelvlakken) en de twee enantiomeren van het cis-isomeer
(de twee chelaten en de twee chloriden winden zich in de ene of de andere richting).
\end{solution}

\begin{solution}{exo:b2:coordination-complexes:11}
Vlakke zeshoek: B op de posities 1,2 (ortho), 1,3 (meta), 1,4 (para): drie isomeren.
Trigonaal prisma: B langs een ribbe van een driehoek, langs een verticale ribbe, of diagonaal over een
rechthoekig zijvlak: drie. De octaëder geeft er twee. Er niet meer dan twee vinden, ondanks
veel pogingen en verscheidene verbindingen, sluit de twee andere geometrieën uit --- tenzij er een
derde isomeer over het hoofd gezien is.
\end{solution}

\begin{solution}{exo:b2:coordination-complexes:12}
\ce{[RhCl(PPh3)3]}: Rh(I) $d^8$ + 4 × 2 = 16. \ce{[IrH(CO)(PPh3)3]}: Ir(I) $d^8$ + 5 × 2 =
18. Het rodiumcomplex, met 16 elektronen, kan er nog twee opnemen (oxidatieve additie,
\cref{ch:b2:catalytic-cycles}).
\end{solution}

\begin{solution}{pb:b2:coordination-complexes:1}
\textbf{1.} De chloriden buiten de \omterm{def:b2:coordination-complexes:coordination-sphere}{coördinatiesfeer}, vrij in oplossing.
\textbf{2.} A: 0; B: 1; C en D: 2.
\textbf{3.} A: $[\ce{Co(NH3)6}]^{3+}$ + 3 \ce{Cl-}: 4 ionen; B: 3; C, D: 2. Dat klopt.
\textbf{4.} +3.
\textbf{5.} $d^6$.
\textbf{6.} Zes in elke verbinding: 6 \ce{NH3}; 5 \ce{NH3} + 1 Cl; 4 \ce{NH3} + 2 Cl.
\textbf{7.} \ce{[Co(NH3)6]Cl3}; \ce{[CoCl(NH3)5]Cl2}; \ce{[CoCl2(NH3)4]Cl} (tweemaal).
\textbf{8.} Hexaamminekobalt(III)chloride; pentaammine-chloridokobalt(III)chloride.
\textbf{9.} Tetraammine-dichloridokobalt(III).
\textbf{10.} Geometrische isomeren (cis-trans).
\textbf{11.} \ce{[CoCl3(NH3)3]}, een neutraal complex: geen ionen, geen chloride dat meteen
neerslaat.
\textbf{12.} Twee: fac en mer.
\textbf{13.} Cis: de twee Cl onder $90^\circ$; trans: onder $180^\circ$.
\textbf{14.} Drie (ortho-, meta- en paraposities op de zeshoek).
\textbf{15.} Drie.
\textbf{16.} Het ondersteunt de octaëder, de enige van de drie die er twee voorspelt.
\textbf{17.} Nee: de afwezigheid van een derde isomeer is negatief bewijs. Een derde isomeer zou
de octaëder uitgesloten hebben.
\textbf{18.} Dat welk het chelaat vormt: een tweetandig ethaandioaat kan slechts twee
cisposities overspannen.
\textbf{19.} Drie en-chelaten rond kobalt, die elk twee cisposities overspannen.
\textbf{20.} Geen van beide: geen spiegelvlak en geen inversiecentrum.
\textbf{21.} Twee enantiomeren, $\Delta$ en $\Lambda$.
\textbf{22.} Dat de zes liganden in drie dimensies geschikt zijn, zoals in een
octaëder; een vlakke schikking zou een achiraal ion geven.
\textbf{23.} Cis: chiraal (twee enantiomeren). Trans: achiraal.
\textbf{24.} Het toonde aan dat optische activiteit een eigenschap is van de moleculaire geometrie in het algemeen,
niet van koolstof, en het bevestigde de octaëder.
\textbf{25.} De octaëder voorspelt $\boldsymbol{2}$ isomeren van \ce{[CoCl2(NH3)4]+}; de
vlakke zeshoek en het trigonale prisma voorspellen er $\boldsymbol{3}$.
\end{solution}
